\begin{table*}[t]
\centering\small
\caption{Official VBPIR core~\cite{vectorbatchpir} with portability changes and an external root-verification adapter, on the common height-10 tree. AB is ours. Times are descriptive means over ten paired targets in one process per layout, using the original C++ timing boundaries and fresh client/server objects per target. No confidence interval is estimated. Estimated query/answer sizes are the original program's serialization-size estimates, not measured wire traffic, and are not directly comparable with SimplePIR matrix bytes or TCP application bytes.}
\label{tab:official-vbpir-components}
\setlength{\tabcolsep}{8pt}
\begin{tabular}{lrrrrr}
\toprule
Layout & Query (ms) & Answer (ms) & Recover (ms) & Est. query (KiB) & Est. answer (KiB)\\
\midrule
TreePIR CSA~\cite{treepir} & 6.761 & 338.488 & 1.263 & 385 & 257\\
AB & 7.206 & 338.839 & 1.211 & 385 & 257\\
\bottomrule
\end{tabular}
\par\smallskip
\begin{minipage}{.97\textwidth}\small
Initialization: N/A, original setup is untimed. Full client state and isolated client RSS: N/A, unmeasured in this combined process. Actual wire bytes and full-service E2E: N/A, no transport and verification is outside the timers. The original estimates use query \texttt{save\_size()/2} and answer \texttt{save\_size()}, then integer division by 1024. All 20 output proofs pass independent root verification and corruption rejection.
\end{minipage}
\end{table*}
